\documentclass[11pt,a4paper]{article}
\usepackage[T1]{fontenc}
\usepackage[utf8]{inputenc}
\usepackage{newtxtext}
\usepackage{amsmath,amsthm,mathtools}
\usepackage{newtxmath}
\usepackage[margin=27mm,headheight=15pt]{geometry}
\usepackage{microtype,booktabs,tabularx,array}
\usepackage[dvipsnames]{xcolor}
\usepackage{enumitem,fancyhdr,hyperref,bookmark}
\hypersetup{colorlinks=true,linkcolor=MidnightBlue,citecolor=MidnightBlue,
 urlcolor=MidnightBlue,pdftitle={Global Fueter Primitives: Period Obstructions, Spherical Defects, and Sharp Energy Laws},
 pdfauthor={OpenAI},pdfsubject={Research extension in quaternionic analysis}}
\urlstyle{same}
\numberwithin{equation}{section}
\newtheorem{theorem}{Theorem}[section]
\newtheorem{proposition}[theorem]{Proposition}
\newtheorem{lemma}[theorem]{Lemma}
\newtheorem{corollary}[theorem]{Corollary}
\theoremstyle{definition}
\newtheorem{definition}[theorem]{Definition}
\newtheorem{example}[theorem]{Example}
\theoremstyle{remark}
\newtheorem{remark}[theorem]{Remark}
\newcommand{\R}{\mathbb R}
\newcommand{\C}{\mathbb C}
\newcommand{\HH}{\mathbb H}
\newcommand{\HC}{\mathbb H_{\C}}
\newcommand{\SSS}{\mathbb S}
\newcommand{\SR}{\mathcal{SR}}
\newcommand{\AM}{\mathcal{AM}}
\newcommand{\D}{\mathcal D}
\newcommand{\lap}{\Delta_4}
\newcommand{\cI}{\mathcal I}
\newcommand{\Per}{\operatorname{Per}}
\newcommand{\wind}{\operatorname{wind}}
\newcommand{\supp}{\operatorname{supp}}
\newcommand{\dd}{\,\mathrm d}
\newcommand{\abs}[1]{\lvert#1\rvert}
\newcommand{\norm}[1]{\lVert#1\rVert}
\newcommand{\qf}{\mathcal Q}
\newcommand{\aff}{\operatorname{Aff}_{\HH}}
\newcommand{\ii}{\mathbf i}
\newcommand{\jj}{\mathbf j}
\newcommand{\kk}{\mathbf k}
\setlength{\parskip}{2pt}
\setlength{\parindent}{1.2em}
\setlength{\emergencystretch}{2em}
\setlist{itemsep=2pt,topsep=5pt}
\pagestyle{fancy}
\fancyhf{}
\fancyhead[L]{\small Global Fueter primitives}
\fancyhead[R]{\small\thepage}
\renewcommand{\headrulewidth}{0.3pt}
\title{\vspace{-10mm}\textbf{Global Fueter Primitives}\\[4pt]
{\Large Period Obstructions, Spherical Defects,\\and Sharp Energy Laws}}
\author{A research extension of the supplied quaternionic-analysis manuscript}
\date{21 September 2026}
\begin{document}
\maketitle
\thispagestyle{empty}
\begin{abstract}
Local inversion of the quaternionic Fueter map does not by itself give a
single-valued global primitive. We identify the obstruction by two explicit,
closed quaternion-valued one-forms. On a conjugation-symmetric planar set of
finite topology, each conjugate pair of complementary holes contributes
exactly two quaternionic obstruction parameters; a hole fixed by conjugation
contributes none. This gives an explicit split description of the range of
the Laplacian, including a necessary and sufficient topological criterion for
surjectivity. Period-normalized spherical kernels realize every obstruction,
and lead to a period-preserving rational approximation theorem.

Near a nonreal singular sphere, we classify all axial Fueter-regular fields
of finite power growth: an affine-logarithmic defect, finitely many rational
Fueter principal parts, and a removable remainder. At the critical growth
order $\rho^{-1}$, the two periods $a,b$ determine the sharp logarithmic
$L^2$-energy coefficient
$8(t^2|a|^2+|sa+b|^2)$ for the sphere $s+t\SSS$. We prove a matching lower
bound without a growth hypothesis, identify the corresponding distributional
surface source, and give an $L^2$ removability argument. All main assertions
are proved; symbolic and numerical diagnostics are supplied separately.
\end{abstract}
\medskip
\noindent\textbf{Scope.} The setting is one quaternionic variable, left slice
regularity with right coefficients, and axial left Cauchy--Fueter regularity.
The energy is the squared $L^2$ norm of the field, not a Dirichlet energy.

\medskip
\noindent\textbf{Research status.} The local inverse and the existence of
sphere-supported Cauchy kernels are established mathematics. The candidate
contributions here are the explicit reflection-sensitive global range
classification, its approximation consequences, the finite-growth normal
form in these period coordinates, and the sharp energy/source formulas.
A targeted literature search did not locate these statements in this combined
form; this is not an exhaustive priority determination. The paper does not
claim to settle a named published conjecture.
\newpage
\begingroup
\small
\tableofcontents
\endgroup
\newpage

\section{The question and its relation to previous work}
\label{sec:intro}
The supplied manuscript \cite{Attachment} develops classical quaternionic
analogues and ends its Fueter section with a constructive inverse on simply
connected subsets of the upper half-plane. It explicitly leaves global
period conditions untreated. We address that gap rather than repeating its
Cauchy formulas, zero theory, or Schwarz--Pick theory.

The basic question is concrete. Given a single-valued axial Fueter-regular
function $g$, when is there a single-valued slice-regular $f$ on the
\emph{same} circular domain such that
\begin{equation}
 \lap f=g?
 \label{eq:problem}
\end{equation}
There is also a quantitative version: how much singular growth or $L^2$
energy is forced when local primitives return with nonzero affine monodromy?
A third question asks whether the resulting defects account for all
obstructions to approximation by Fueter images of rational slice functions.

\subsection{Established ingredients and the proposed increment}
The inverse Fueter mapping theorem was developed by Colombo, Sabadini and
Sommen \cite{CSS11,CSS13}. Colombo, Pe\~na Pe\~na, Sabadini and Sommen
\cite{CPSS14} gave radial integral formulas, described the kernel, and
computed primitives of kernels singular on a sphere. Their preprint develops
its explicit construction on rectangles. These facts, and not a claim that
the inverse transform itself is new, form our starting point.

The closest elementary global result located in the search is Perotti's
Theorem~2 in \cite{Perotti24}: surjectivity holds when every connected
component of the conjugation-symmetric planar set is simply connected. Its
proof constructs a potential and then a harmonic conjugate. We replace
that sufficient condition by necessary and sufficient period conditions,
and, in finite topology, determine exactly which holes matter.

The stem framework used below is the one developed in \cite{GP11}.
Recent work of De Martino and Pinton \cite{DMP26}, posted in August 2026,
studies factorizations of the Fueter map and order-two polyanalytic
extensions. Our operator remains the ordinary four-dimensional Laplacian;
we do not claim results for those factorized or higher-order theories.

There are two important qualifications. First, the passage from closed
forms to periods is standard topology, and the rational approximation step
uses the classical Runge theorem. Second, the logarithmic kernels below
are closely related to the already studied spherical Cauchy kernels.
The proposed increment is an explicit range theorem with reflection, a
complete normalization of its defects, and the precise singularity and
energy consequences. The references and the accompanying literature log
make this distinction explicit.

\subsection{A preview of the main results}
Write an axial field as $g(x+Iy)=P(x,y)+I Q(x,y)$ and put
\begin{align}
 \omega_g&=\frac12(-P\dd x+Q\dd y),\\
 \theta_g&=\frac12\big((yQ+xP)\dd x+(yP-xQ)\dd y\big).
 \label{eq:previewforms}
\end{align}
The two forms are closed and invariant under reflection $y\mapsto-y$.
A global primitive exists exactly when both have zero periods. For a
finite-type planar set, let $h$ be the number of conjugate pairs of holes.
We will prove
\begin{equation}
 \AM(\Omega_D)/\lap\SR(\Omega_D)\simeq\HH^{2h}.
 \label{eq:previewquotient}
\end{equation}
Thus an ordinary round annulus centered on the real axis has no obstruction,
but deleting one nonreal conjugate pair does produce an obstruction.

At a puncture $\zeta=s+\iota t$ in the upper half-plane, corresponding to
the sphere $s+t\SSS$, let $a,b$ be the two periods around a positively
oriented small circle. If $\rho=|z-\zeta|$ and $g=O(\rho^{-1})$, then
\begin{equation}
 \int_{\{\varepsilon<\rho<R\}}|g(q)|^2\dd V(q)
 =8\qf_\zeta(a,b)\log(R/\varepsilon)+C_g(R)+o(1),
 \qquad
 \qf_\zeta(a,b)=t^2|a|^2+|sa+b|^2.
 \label{eq:previewenergy}
\end{equation}
The integration set in this formula is the four-dimensional
circularization of the planar annulus. The coefficient is positive for
every nonzero period pair.

\section{Conventions, stems, and axial differential identities}
\label{sec:prelim}
The quaternionic units satisfy $\ii^2=\jj^2=\kk^2=\ii\jj\kk=-1$.
We use $\iota$ for the imaginary unit in a separate commutative complex
field, and write
\[
 \HC=\HH\otimes_\R\C=\{A+\iota B:A,B\in\HH\}.
\]
Complex conjugation on $\HC$ is
$\overline{A+\iota B}=A-\iota B$; it is \emph{not} quaternionic
conjugation. For estimates only, use its Euclidean norm
$\|A+\iota B\|^2=|A|^2+|B|^2$. Multiplication by a complex number of
modulus one preserves this norm.

Let $D\subset\C$ be open and invariant under
$\tau(z)=\bar z$. Its circularization is
\[
 \Omega_D=\{x+Iy:x+\iota y\in D,\ I\in\SSS\},
 \qquad \SSS=\{I\in\HH:I^2=-1\}.
\]
Except where a local statement says otherwise, $\Omega_D$ is connected.
Then $D$ is either connected and meets $\R$, or consists of two conjugate
connected components, one in each half-plane. The latter case is a
product domain.

A stem $F=A+\iota B:D\to\HC$ obeys
\[
 A(x,-y)=A(x,y),\qquad B(x,-y)=-B(x,y).
\]
It induces $f=\cI(F)$ by $f(x+Iy)=A(x,y)+I B(x,y)$.
It is holomorphic precisely when
\begin{equation}
 A_x=B_y,\qquad A_y=-B_x.
 \label{eq:CR}
\end{equation}
We define $\SR(\Omega_D)$ to be the right $\HH$-module of functions
induced by holomorphic stems. On a domain meeting $\R$ this agrees with
the usual slice-wise definition; on a product domain the stem condition
is part of the definition.

The left Cauchy--Fueter operator and its conjugate are
\[
 \D=\partial_x+\ii\partial_{x_1}+\jj\partial_{x_2}+\kk\partial_{x_3},
 \qquad
 \overline{\D}=\partial_x-\ii\partial_{x_1}-\jj\partial_{x_2}-\kk\partial_{x_3},
 \qquad \overline{\D}\D=\lap.
\]
Let $\AM(\Omega_D)$ denote the smooth slice functions $g=P+I Q$ with
$\D g=0$. Here $P,Q$ are quaternion-valued, not necessarily real-valued.
At the real axis their parity is understood through the inducing smooth
stem. Coefficients and all period parameters are always multiplied on the
right.

\begin{lemma}[Axial calculus]
\label{lem:axial}
For $y>0$ and a smooth axial function $u=A+I B$,
\begin{align}
 \D u&=\left(A_x-B_y-\frac{2B}{y}\right)+I(B_x+A_y),\\
 \lap u&=A_{xx}+A_{yy}+\frac{2A_y}{y}
 +I\left(B_{xx}+B_{yy}+\frac{2B_y}{y}-\frac{2B}{y^2}\right).
 \label{eq:axial}
\end{align}
Consequently $g=P+I Q$ is Fueter regular exactly when
\begin{equation}
 P_x-Q_y-\frac{2Q}{y}=0,\qquad Q_x+P_y=0,
 \label{eq:Vekua}
\end{equation}
and a slice-regular $f=A+IB$ satisfies
\begin{equation}
 \lap f=\frac{2A_y}{y}
 +I\left(\frac{2B_y}{y}-\frac{2B}{y^2}\right).
 \label{eq:Fueter}
\end{equation}
The last expression is Fueter regular.
\end{lemma}
\begin{proof}
Put $v=\ii x_1+\jj x_2+\kk x_3$, $y=|v|$, and $I=v/y$.
For $\ell=1,2,3$,
\[
 \partial_{x_\ell}y=x_\ell/y,\qquad
 \partial_{x_\ell}I=e_\ell/y-vx_\ell/y^3.
\]
Summation gives $\sum e_\ell\partial_{x_\ell}I=-2/y$,
$\Delta_{\R^3}I=-2I/y^2$, and
$\sum(\partial_{x_\ell}I)(\partial_{x_\ell}y)=0$.
The product rule proves \eqref{eq:axial}. Comparing the values at $I$ and
$-I$ proves \eqref{eq:Vekua}. For a holomorphic stem both $A$ and $B$
are planar harmonic, giving \eqref{eq:Fueter}. Substitution of its two
coefficients into \eqref{eq:Vekua}, using \eqref{eq:CR}, gives zero.
Where the domain meets $\R$, smoothness extends the identities across
$y=0$ by continuity. The formulas with signed $y\ne0$ follow by parity.
\end{proof}

\begin{lemma}[Spherical norm identity]
\label{lem:norm}
For $P,Q\in\HH$,
\begin{equation}
 \frac1{4\pi}\int_{\SSS}|P+IQ|^2\dd\sigma(I)=|P|^2+|Q|^2,
 \label{eq:meansquare}
\end{equation}
where $\sigma(\SSS)=4\pi$. Hence
\[
 (|P|^2+|Q|^2)^{1/2}
 \leq\sup_{I\in\SSS}|P+IQ|
 \leq\sqrt2(|P|^2+|Q|^2)^{1/2}.
\]
\end{lemma}
\begin{proof}
Expand the square. The mixed real inner product is linear in $I$ and has
zero average since $\int_{\SSS}I\dd\sigma=0$; also $|IQ|=|Q|$.
The upper bound is the triangle inequality followed by Cauchy--Schwarz.
\end{proof}

\section{Two closed forms and the complete global inverse criterion}
\label{sec:periods}
Define $\omega_g,\theta_g$ by \eqref{eq:previewforms} on all of $D$,
using signed $y$. Unlike the Vekua equations, these expressions have no
division by $y$.

\begin{lemma}[Closedness and reflection]
\label{lem:closed}
For $g\in\AM(\Omega_D)$,
\[
 d\omega_g=d\theta_g=0,\qquad
 \tau^*\omega_g=\omega_g,\qquad\tau^*\theta_g=\theta_g.
\]
If $g=\lap f$ and $f=\cI(A+\iota B)$, then the smooth even functions
\begin{equation}
 C=B/y,\qquad V=A-xB/y
 \label{eq:CV}
\end{equation}
satisfy $dC=\omega_g$ and $dV=\theta_g$. At $y=0$, $B/y$ is interpreted
as $B_y(x,0)$.
\end{lemma}
\begin{proof}
Writing both exterior derivatives with orientation $dx\wedge dy$ gives
\begin{align*}
 d\omega_g&=\tfrac12(Q_x+P_y)\,dx\wedge dy,\\
 d\theta_g&=\tfrac12\big[y(P_x-Q_y)-2Q-x(Q_x+P_y)\big]dx\wedge dy.
\end{align*}
They vanish by \eqref{eq:Vekua} off the real axis, hence everywhere.
The parity of $P,Q$ proves reflection invariance, including the sign
change of $dy$.

For a Fueter image, \eqref{eq:Fueter} and \eqref{eq:CR} yield
\[
 C_x=-P/2,\quad C_y=Q/2,\quad
 V_x=(yQ+xP)/2,\quad V_y=(yP-xQ)/2.
\]
The extension of $C$ is smooth because locally
$B(x,y)=y\int_0^1 B_y(x,ty)\dd t$. Its parity, and that of $V$, is even.
\end{proof}

\begin{theorem}[Global period criterion and explicit reconstruction]
\label{thm:global}
For $g\in\AM(\Omega_D)$ the following statements are equivalent:
\begin{enumerate}[label=(\roman*)]
 \item There exists $f\in\SR(\Omega_D)$ with $\lap f=g$.
 \item Both $\omega_g$ and $\theta_g$ are exact on every component of $D$.
 \item For every piecewise smooth closed curve $\gamma\subset D$,
 \[
 a_\gamma(g):=\int_\gamma\omega_g=0,
 \qquad b_\gamma(g):=\int_\gamma\theta_g=0.
 \]
\end{enumerate}
When these conditions hold, choose even primitives $C,V$ of the two forms.
Then every solution is of the form
\begin{equation}
 f(x+Iy)=(x+Iy)C(x,y)+V(x,y)+qa+b,
 \qquad a,b\in\HH,
 \label{eq:reconstruction}
\end{equation}
where $q=x+Iy$. In particular,
$\ker(\lap|_{\SR})=\aff=\{q\mapsto qa+b:a,b\in\HH\}$.
\end{theorem}
\begin{proof}
Lemma~\ref{lem:closed} gives (i)$\Rightarrow$(ii). Exactness implies zero
periods. Conversely, fix a basepoint in each component and integrate the
closed form along a path to the endpoint. The difference between two
paths is a closed curve, so (iii) makes the integral path independent;
local differentiation gives a primitive. This proves (iii)$\Rightarrow$(ii).

If $C_0,V_0$ are any primitives on $D$, replace each by its average with
its pullback under $\tau$. Reflection invariance of the forms ensures
that $C=(C_0+C_0\circ\tau)/2$ and similarly $V$ are still primitives,
and are even. Set $A=xC+V$ and $B=yC$. Their derivatives are
\[
 A_x=C+yQ/2=B_y,\qquad A_y=yP/2=-B_x.
\]
Thus $A+\iota B$ is a holomorphic stem, even across the real axis.
Equation~\eqref{eq:Fueter} gives $\lap f=g$ off that axis, and then on it
by continuity.

If $\lap f=0$, the functions in \eqref{eq:CV} have zero differentials.
They are constant on each component; evenness identifies their constants
on conjugate components. Since $\Omega_D$ is connected, this yields
$f(q)=qa+b$. Conversely every such function has zero Laplacian.
\end{proof}

\begin{remark}[Why there are two independent periods]
The supplied local proof first integrates $\omega_g$ and then integrates a
form containing the first primitive. Passing from $A$ to $V=A-xB/y$
removes that dependence: both forms above are computed directly from $g$.
This is not merely a change in notation. It exposes two independent
period obstructions before any path integration has been performed.
The kernels in Section~\ref{sec:kernels} realize $(a,b)=(1,0)$ and
$(0,1)$ separately.
\end{remark}

\begin{corollary}[Affine monodromy]
\label{cor:monodromy}
Continue a local Fueter primitive along a closed curve $\gamma$ in the
planar parameter domain. On returning, its stem changes by
$z a_\gamma(g)+b_\gamma(g)$, and its induced slice function changes by
$qa_\gamma(g)+b_\gamma(g)$. Both period maps factor through
$H_1(D;\mathbb Z)$ and satisfy
\begin{equation}
 a_{\tau\gamma}(g)=a_\gamma(g),\qquad
 b_{\tau\gamma}(g)=b_\gamma(g).
 \label{eq:reflectionperiod}
\end{equation}
\end{corollary}
\begin{proof}
Continue the path-integral primitives $C,V$ in
\eqref{eq:reconstruction}. Their changes are the respective integrals.
Stokes' theorem for the closed forms gives homology invariance, and
pullback by $\tau$ gives \eqref{eq:reflectionperiod}.
\end{proof}

\subsection{Quantitative continuity of the inverse}
For a compact planar set $K$, write
$\|g\|_{\Omega_K}=\sup\{|g(x+Iy)|:x+\iota y\in K,I\in\SSS\}$.
Suppose normalized primitives are obtained using paths in a compact
$K'\subset D$ from a basepoint to every point of $K$, with lengths at
most $L$. Write $M=\|g\|_{\Omega_{K'}}$ and
$R=\max_{K'}|z|$. The norms of the coefficient pairs of $\omega_g$ and
$\theta_g$ are at most $M/2$ and $RM/2$, respectively. Therefore
\begin{equation}
 |C|\leq LM/2,
 \qquad |V|\leq LRM/2,
 \qquad \|f\|_{\Omega_K}\leq LRM.
 \label{eq:inversebound}
\end{equation}
For any fixed compact $K$, such a larger compact $K'$ and a finite $L$
exist: cover $K$ by finitely many small disks and join their centers to
the basepoint by finitely many paths. In a connected symmetric $D$ one
may take a real basepoint; the normalization is then
$f(x_*)=\partial_c f(x_*)=0$. In a product domain, normalize $C,V$ at
one upper basepoint and reflect. Thus the inverse modulo affine functions
is continuous on compact sets; derivatives are controlled by the displayed
differential identities and local holomorphic estimates.

\section{Period-normalized spherical kernels}
\label{sec:kernels}
We first work in $y>0$. Let $\zeta=s+\iota t$, initially with $t>0$,
and let $a,b\in\HH$. A local branch of the holomorphic function
\begin{equation}
 F_\zeta^{a,b}(z)=\frac{za+b}{2\pi\iota}\log(z-\zeta)
 \label{eq:logstem}
\end{equation}
changes by $za+b$ when continued once counterclockwise around $\zeta$.
It follows from the affine kernel that the local Fueter images fit
together to give a single-valued field $G_\zeta^{a,b}$ on the
circularization of $\C^+\setminus\{\zeta\}$.

\begin{proposition}[Branch-free formula]
\label{prop:kernel}
Put $\xi=x-s$, $\eta=y-t$, $\rho^2=\xi^2+\eta^2$, $X=xa+b$,
and $\ell=\log\rho$. Then $G_\zeta^{a,b}=P_\zeta^{a,b}+I Q_\zeta^{a,b}$,
where
\begin{align}
 P_\zeta^{a,b}
 &=\frac{X\xi+ya\eta}{\pi y\rho^2}+\frac{a\ell}{\pi y},\\
 Q_\zeta^{a,b}
 &=\frac{ya\xi-X\eta}{\pi y\rho^2}+\frac{X\ell}{\pi y^2}.
 \label{eq:kernelPQ}
\end{align}
For every upper-half-plane closed curve avoiding $\zeta$,
\begin{equation}
 \big(a_\gamma(G_\zeta^{a,b}),b_\gamma(G_\zeta^{a,b})\big)
 =\wind(\gamma,\zeta)(a,b).
 \label{eq:kernelperiod}
\end{equation}
Formula~\eqref{eq:kernelPQ} also defines $G_\zeta^{a,b}$ for $t<0$ on
all of $\C^+$; in that case it has zero periods there.
\end{proposition}
\begin{proof}
Write $\log(z-\zeta)=\ell+\iota v$ on a local branch. The real and
imaginary stem components in \eqref{eq:logstem} are
\[
 A=\frac{Xv+ya\ell}{2\pi},\qquad
 B=\frac{yav-X\ell}{2\pi}.
\]
Since $v_y=\xi/\rho^2$ and $\ell_y=\eta/\rho^2$,
\eqref{eq:Fueter} gives \eqref{eq:kernelPQ}; the argument $v$ cancels.
The period formula follows either from continuation and
Corollary~\ref{cor:monodromy}, or from the differentials of $B/y$ and
$A-xB/y$. If $t<0$, a branch of $\log(z-\zeta)$ is holomorphic throughout
$\C^+$, so the two periods vanish.
\end{proof}

The bare kernel need not extend smoothly across the real axis. The
reflected difference does.

\begin{theorem}[A global spherical defect]
\label{thm:pairedkernel}
For $t>0$, the field
\begin{equation}
 K_\zeta^{a,b}=G_\zeta^{a,b}-G_{\bar\zeta}^{a,b}
 \label{eq:pairedkernel}
\end{equation}
on $y>0$ extends to an axial Fueter-regular function on
$\HH\setminus\Sigma_\zeta$, where
$\Sigma_\zeta=s+t\SSS$. It is locally the Fueter image of
\begin{equation}
 H_\zeta^{a,b}(z)=\frac{za+b}{2\pi\iota}
 \log\frac{z-\zeta}{z-\bar\zeta}.
 \label{eq:pairedlog}
\end{equation}
Its periods around $\zeta$ are $(a,b)$ and around $\bar\zeta$ are
$(-a,-b)$, for counterclockwise circles in both cases.
\end{theorem}
\begin{proof}
Away from the real axis, subtract the two local logarithms and apply
Proposition~\ref{prop:kernel}. Near a real point, the quotient
$R(z)=(z-\zeta)/(z-\bar\zeta)$ is nonzero and satisfies
$R(\bar z)=1/\overline{R(z)}$. Choose a local logarithm $L$ with a purely
imaginary value at that real point. Uniqueness of the branch then gives
$L(\bar z)=-\overline{L(z)}$. Hence \eqref{eq:pairedlog} is a
holomorphic stem on that neighborhood. Different choices of logarithm
change it by an integer multiple of $za+b$, whose Fueter image is zero.
These local images therefore patch smoothly over the real axis and
across all other overlaps. The period statements follow from the winding
numbers of the numerator and denominator of $R$.
\end{proof}

\begin{example}[A counterexample on a genuine slice domain]
\label{ex:counter}
Take $\zeta=\iota$, $a=0$, $b=1$. On $y>0$, set
$d_-^2=x^2+(y-1)^2$ and $d_+^2=x^2+(y+1)^2$. The field $K=P+I Q$ is
\begin{align}
 P&=\frac{x}{\pi y}\left(\frac1{d_-^2}-\frac1{d_+^2}\right),\\
 Q&=\frac1{\pi y}\left(-\frac{y-1}{d_-^2}
                         +\frac{y+1}{d_+^2}\right)
       +\frac1{\pi y^2}\log\frac{d_-}{d_+}.
 \label{eq:explicitcounter}
\end{align}
It extends smoothly to $\HH\setminus\SSS$, with real-axis value
\[
 K(x)=\frac{4x}{\pi(x^2+1)^2}.
\]
Nevertheless it has no single-valued slice-regular primitive on that
same domain, since its second period is $1$. Its first period is $0$.
The field $K_\iota^{1,0}$ conversely has first period $1$ and second
period $0$. Thus neither of the two tests can be discarded.
\end{example}

\begin{remark}[These kernels are not claimed to be new functions]
The kernels singular on a sphere and their logarithmic primitives already
appear in inverse Fueter theory \cite{CPSS14}. Here the point is their
normalization by \emph{two global periods}, their reflection-compatible
extension, and the role they play as a basis of a cokernel. Section~\ref{sec:source}
identifies them directly as spherical Cauchy layers.
\end{remark}

\section{Finite topology: which holes actually obstruct inversion?}
\label{sec:topology}
For a completely explicit topological statement, we use the following
standard finite-type class. The complement of $D$ in the Riemann sphere
has finitely many components; each bounded complementary component is a
closed Jordan region or a point, and the boundary curves are piecewise
$C^2$. The unbounded complementary component may reduce to the point at
infinity. The same convention applies to a conjugate pair of components
of $D$. One may equivalently work with finite smooth planar domains with
finitely many punctures. No assertion about wild boundaries is needed.

A \emph{hole} is a complementary component other than the component
containing infinity. Conjugation either fixes a hole or interchanges it
with a distinct hole. Let $h=h(D)$ be the number of interchanged pairs.
Choose one point $\zeta_j\in\C^+$ in each upper member, and a positively
oriented loop $\gamma_j\subset D$ surrounding that hole and no other.

\begin{lemma}[Reflection on planar homology]
\label{lem:homology}
The classes of positively oriented loops around the holes freely generate
$H_1(D;\mathbb Z)$. If $e_H$ is the class of such a loop around $H$,
then
\[
 \tau_*e_H=-e_{\tau H}.
\]
Consequently an invariant closed form has zero period around a fixed
hole, and opposite periods around a conjugate pair.
\end{lemma}
\begin{proof}
For a finite planar domain, cut along disjoint arcs joining the inner
boundary components and punctures to an outer boundary, or to a large
circle in the unbounded case. The resulting pieces are simply connected.
Restoring the arcs gives the usual independent generators around the
holes. This is also characterized by the integer winding numbers about
one selected point in each complementary component. Conjugation changes
counterclockwise orientation to clockwise orientation and interchanges
the selected holes. The final statements follow from
$\int_\gamma\alpha=\int_{\tau\gamma}\alpha$ and the absence of
$2$-torsion in the additive quaternionic group.
\end{proof}

\begin{theorem}[Complete finite-topology range theorem]
\label{thm:range}
Under the finite-type hypotheses above, the period map
\[
 \Per(g)=\big(a_{\gamma_1}(g),b_{\gamma_1}(g),\ldots,
                    a_{\gamma_h}(g),b_{\gamma_h}(g)\big)
\]
gives an exact sequence of right quaternionic vector spaces
\begin{equation}
 0\longrightarrow\HH^2\longrightarrow\SR(\Omega_D)
 \xrightarrow{\ \lap\ }\AM(\Omega_D)
 \xrightarrow{\ \Per\ }\HH^{2h}\longrightarrow0.
 \label{eq:exact}
\end{equation}
The first map sends $(a,b)$ to $qa+b$. An explicit right inverse to
$\Per$ is
\begin{equation}
 J((a_j,b_j)_{j=1}^h)=\sum_{j=1}^h K_{\zeta_j}^{a_j,b_j}.
 \label{eq:splitting}
\end{equation}
In particular, every $g\in\AM(\Omega_D)$ has a decomposition
\begin{equation}
 g=\lap f+\sum_{j=1}^h K_{\zeta_j}^{a_j,b_j},
 \label{eq:decomposition}
\end{equation}
where $(a_j,b_j)$ are uniquely determined, and $f$ is unique modulo
$\aff$. The cokernel has quaternionic dimension $2h$ and real
dimension $8h$.
\end{theorem}
\begin{proof}
The kernel at $\SR$ is Theorem~\ref{thm:global}. Lemmas
\ref{lem:closed} and \ref{lem:homology} say that a field with zero
selected upper periods has zero periods on every generator, including
fixed holes and lower holes. The global criterion thus identifies
$\ker\Per$ with $\lap\SR$.

The field $K_{\zeta_j}^{a,b}$ is regular on $\Omega_D$ because its
singular sphere lies over a missing complementary component. Its selected
period pair is $(a,b)$ at the $j$th hole and zero at all other holes.
Theorem~\ref{thm:pairedkernel} therefore gives $\Per J=\mathrm{id}$.
Subtracting $J\Per g$ leaves zero periods and proves the decomposition.
Uniqueness of its coefficients and the affine ambiguity follow in the
same order. Right linearity is immediate from the displayed formulas.
\end{proof}

\begin{corollary}[Necessary and sufficient topological condition]
\label{cor:onto}
For a domain in this finite-type class, the Fueter map is onto
$\AM(\Omega_D)$ if and only if every complementary component of $D$
is fixed by complex conjugation. Simple connectivity is sufficient but
not necessary.
\end{corollary}
\begin{proof}
This is equivalent to $h(D)=0$. Apply Theorem~\ref{thm:range}; for the
necessity direction a nonzero kernel field in any missing pair is an
explicit counterexample to surjectivity.
\end{proof}

\begin{example}[Annuli versus spherical punctures]
For $D=\{z:r_0<|z|<r_1\}$, the single bounded hole is invariant under
conjugation. Every axial Fueter-regular function on the quaternionic
annulus has a global slice-regular primitive. In contrast,
$D=\C\setminus\{\iota,-\iota\}$ has one interchanged pair, so its
cokernel is $\HH^2$. Deleting finitely many real points produces no
obstruction of this kind; deleting $h$ distinct nonreal conjugate pairs
produces $2h$ quaternionic parameters.

For a product domain generated by a planar $U\subset\C^+$ with $m$
holes, $D=U\cup\bar U$ has $h=m$. Thus the product-domain cokernel is
$\HH^{2m}$, as the upper-half-plane period calculation suggests.
\end{example}

\begin{remark}[Cohomological formulation]
Let $H^1_{\mathrm{dR}}(D;\R)^+$ denote the $+1$ eigenspace of reflection
on de Rham cohomology. In the finite-type setting, the pair
$([\omega_g],[\theta_g])$ identifies the cokernel with
\[
 H^1_{\mathrm{dR}}(D;\R)^+\otimes_\R\HH^2.
\]
This is a restatement of the proved period theorem, not an invocation of
an unproved quaternionic sheaf theorem. Its dimension is $8h$ because
reflection acts with a minus sign on a fixed hole and as minus-interchange
on each pair. Formula~\eqref{eq:splitting} gives actual representatives.
\end{remark}

\section{Closed range and period-preserving rational approximation}
\label{sec:approx}
\begin{proposition}[A stable obstruction to approximation]
\label{prop:closedrange}
The range $\lap\SR(\Omega_D)$ is closed in $\AM(\Omega_D)$ for the
compact-open topology, and hence also for the smooth compact-open
topology. If $\gamma$ has length $L>0$ and
$R_\gamma=\max_\gamma|z|>0$, then every $f\in\SR(\Omega_D)$ satisfies
\begin{equation}
 \|g-\lap f\|_{\Omega_\gamma}
 \geq\max\left\{\frac{2|a_\gamma(g)|}{L},
           \frac{2|b_\gamma(g)|}{LR_\gamma}\right\}.
 \label{eq:nondensity}
\end{equation}
\end{proposition}
\begin{proof}
Let $u=P+I Q$ be an axial field and put
$M=\|u\|_{\Omega_\gamma}$. Lemma~\ref{lem:norm} gives
$(|P|^2+|Q|^2)^{1/2}\leq M$. The coefficient-pair norms of its two
forms are at most $M/2$ and $|z|M/2$. Integrating yields
\[
 |a_\gamma(u)|\leq LM/2,\qquad
 |b_\gamma(u)|\leq LR_\gamma M/2.
\]
Apply this to $u=g-\lap f$; the periods of $\lap f$ vanish. Each period
functional is continuous for compact-open convergence, so their common
kernel is closed. Theorem~\ref{thm:global} identifies that common kernel
with the range, even without finite topology.
\end{proof}

Choose a conjugation-invariant set $A$ of allowed poles, containing one
real point in every invariant hole and the two points
$\zeta_j,\bar\zeta_j$ for every interchanged pair; allow a pole at
infinity as well. Every invariant hole meets the real axis: a connected
reflection-invariant set avoiding that axis would split into disjoint upper
and lower parts. Thus the required real point exists. A \emph{rational stem} is a rational $\HC$-valued
function $R$ with $R(\bar z)=\overline{R(z)}$. Its induced function
$r=\cI(R)$ is a rational slice function, regular on $\Omega_D$ when its
poles lie in $A$.

\begin{theorem}[Corrected Runge theorem]
\label{thm:runge}
For every $g\in\AM(\Omega_D)$ in the finite-type setting there are
rational stems $R_n$, with poles only in the chosen set $A$, such that
\begin{equation}
 g_n=\lap\cI(R_n)+J\Per(g)\longrightarrow g
 \label{eq:runge}
\end{equation}
with all derivatives locally uniformly on $\Omega_D$. Every $g_n$ has
exactly the same period vector as $g$. Approximation by the terms
$\lap\cI(R_n)$ alone is possible if and only if $\Per(g)=0$.
The $2h$ fixed quaternionic kernel directions used by $J$ are the minimum
number of fixed right-quaternionic directions that can correct all
period vectors.
\end{theorem}
\begin{proof}
By Theorem~\ref{thm:range}, $g-J\Per(g)=\lap f$ for a holomorphic stem
$F$ inducing $f$. Apply the classical Runge approximation theorem to the
four complex coordinate functions of $F$, simultaneously on a symmetric
compact exhaustion of $D$, using the chosen pole in each complementary
component. It gives rational $S_n\to F$ locally uniformly. Replace
\[
 R_n(z)=\tfrac12\big(S_n(z)+\overline{S_n(\bar z)}\big).
\]
The pole set is preserved, $R_n$ is a stem, and $R_n\to F$.
Cauchy's estimates give convergence of all holomorphic derivatives.
Away from the real axis the axial formulas give convergence of all
real derivatives of the induced functions. Near a real point, use the
convergent real-centered Taylor expansion of the stem; termwise
quaternionic differentiation and Cauchy's coefficient bounds give the
same conclusion. A finite cover handles any compact set. Applying
$\lap$ proves \eqref{eq:runge}.

The period assertion follows from $\Per\lap=0$. Necessity for uncorrected
approximation follows from Proposition~\ref{prop:closedrange}; sufficiency
was just proved with $J\Per(g)=0$. Finally any fixed correcting family
must span the cokernel $\HH^{2h}$ under $\Per$, so it has at least $2h$
right-quaternionic directions. The displayed family attains that number.
\end{proof}

\begin{remark}
The Runge theorem is a classical input here, not a newly proved complex
approximation theorem. The content is the exact, finite-dimensional
correction required after applying the Fueter map. A small numerical
least-squares residual cannot bypass a nonzero period: inequality
\eqref{eq:nondensity} gives a nonzero error floor on any detecting loop.
\end{remark}

\section{Finite-order singularities along a nonreal sphere}
\label{sec:singular}
Fix $\zeta=s+\iota t$ with $t>0$ and $0<R<t$. Write
\[
 B_R^*=\{z:0<|z-\zeta|<R\},\qquad
 T_{\varepsilon,R}=\{x+Iy:\varepsilon<|x+\iota y-\zeta|<R\}.
\]
The parameter $\rho=|z-\zeta|$ is exactly the Euclidean distance from
$q=x+Iy$ to $\Sigma_\zeta=s+t\SSS$. Every growth estimate in this
section is uniform in the planar angle and in $I$.

For $C\in\HC$ and $k\geq1$, introduce the symmetric rational stem
\begin{equation}
 R_{\zeta,k,C}(z)=\frac{C}{(z-\zeta)^k}
                 +\frac{\bar C}{(z-\bar\zeta)^k},
 \qquad
 L_{\zeta,k,C}=\lap\cI(R_{\zeta,k,C}).
 \label{eq:rationalprincipal}
\end{equation}
The bar on $C$ is complex-stem conjugation only. The second summand is
regular near $\zeta$ and ensures the correct reflection globally.

\begin{lemma}[Growth of a zero-period primitive]
\label{lem:primitivegrowth}
Suppose $u\in\AM(\Omega_{B_R^*\cup\bar B_R^*})$ has zero periods and
$u=O(\rho^{-N})$ for an integer $N\geq1$. There is a single-valued
holomorphic upper stem $F$ with $\lap\cI(F)=u$ such that
\begin{equation}
 \|F(z)\|=
 \begin{cases}
 O(1+|\log\rho|),&N=1,\\
 O(\rho^{1-N}),&N\geq2.
 \end{cases}
 \label{eq:primitivegrowth}
\end{equation}
Consequently $F$ has a removable singularity for $N=1$, and has a pole
of order at most $N-1$ for $N\geq2$.
\end{lemma}
\begin{proof}
By the global criterion choose single-valued primitives $C,V$ for the
two forms on the punctured disk. Here $y$ is bounded below by $t-R>0$
and $|z|$ is bounded, so their differential norms are $O(\rho^{-N})$.
From a fixed reference circle follow a fixed ray to radius $\rho$ and
then an arc of length at most $2\pi\rho$ to the desired point. The
radial integral is $O(1+|\log\rho|)$ for $N=1$, and
$O(\rho^{1-N})$ for $N\geq2$; the final arc contributes
$O(\rho^{1-N})$. Thus $F=zC+V$ has \eqref{eq:primitivegrowth}.

For each complex coordinate, a negative Laurent coefficient of order
$k$ is bounded in absolute value by a constant times
$\rho^k\sup_{|z-\zeta|=\rho}\|F(z)\|$. For $N=1$ this tends to zero
for every $k\geq1$. For $N\geq2$ it tends to zero for $k>N-1$.
This proves the two assertions without assuming meromorphicity in advance.
\end{proof}

\begin{theorem}[Complete finite-growth normal form]
\label{thm:normalform}
Let $N\geq1$ be an integer and let $g$ be axial Fueter regular on the
punctured tube, with $g=O(\rho^{-N})$. Let $(a,b)$ be its period pair
around a counterclockwise circle about $\zeta$. There are uniquely
determined $C_1,\ldots,C_{N-1}\in\HC$ and a uniquely determined field
$h$, axial Fueter regular on the full tube $\rho<R$, such that
\begin{equation}
 g=K_\zeta^{a,b}+\sum_{k=1}^{N-1}L_{\zeta,k,C_k}+h.
 \label{eq:normalform}
\end{equation}
For $N=1$ the sum is empty. Conversely every such expression has growth
$O(\rho^{-N})$. Modulo removable fields, the space of growth order at
most $N$ has right-quaternionic dimension $2N$ and real dimension $8N$.
\end{theorem}
\begin{proof}
The explicit kernel satisfies $K_\zeta^{a,b}=O(\rho^{-1})$, since its
logarithmic terms are smaller than $\rho^{-1}$. Subtract it. The remaining
field $u$ has zero periods and growth $O(\rho^{-N})$. By
Lemma~\ref{lem:primitivegrowth} it has a holomorphic primitive whose
principal part is
$\sum_{k=1}^{N-1}(z-\zeta)^{-k}C_k$. Subtract the symmetric rational
stems in \eqref{eq:rationalprincipal}. Their reflected summands are
holomorphic in the disk about $\zeta$, so the resulting stem is
holomorphic throughout that disk. Its Fueter image is the required $h$.

The periods uniquely determine $a,b$. If two choices of the $C_k$ gave
removable remainders, subtract the choices. The removable difference
field has a holomorphic Fueter primitive on the full disk by
Theorem~\ref{thm:global}. Subtract that primitive from the rational
principal part. Its Fueter image is zero on the punctured disk, so it is
affine. An affine stem has no negative Laurent coefficients; therefore
all coefficient differences vanish. The remainder field is then fixed.

A pole of order $k$ in a stem gives, by \eqref{eq:Fueter}, growth at most
$\rho^{-k-1}$. This proves the converse. There are two quaternionic
parameters $a,b$ and two for each $C_k\in\HC=\HH\oplus\iota\HH$.
Existence and uniqueness show that all are independent modulo removable
fields, giving $2+2(N-1)=2N$.
\end{proof}

\begin{corollary}[Critical-growth classification]
\label{cor:critical}
Every axial Fueter-regular field with $g=O(\rho^{-1})$ is a uniquely
specified spherical defect $K_\zeta^{a,b}$ plus a removable field.
Such a field is removable if and only if $a=b=0$. In particular,
$g=o(\rho^{-1})$ is removable, and the little-oh threshold is sharp.
\end{corollary}
\begin{proof}
Apply Theorem~\ref{thm:normalform} with $N=1$. For the little-oh claim,
the lengths of the shrinking circles are $2\pi\rho$ and the coefficients
of the two forms are bounded by a fixed multiple of $|g|$ in spherical
supremum norm. Both periods tend to zero and, being independent of
radius, are zero. A nonzero $K_\zeta^{a,b}$ is a sharp counterexample
with big-oh growth.
\end{proof}

\section{Sharp energy forced by monodromy}
\label{sec:energy}
Throughout this section the field is defined on $0<\rho<R_0<t$, and
any outer energy radius $R$ is chosen strictly smaller than $R_0$.
The four-dimensional volume form in the axial coordinates is
$\dd V=y^2\dd x\dd y\dd\sigma(I)$. Lemma~\ref{lem:norm} therefore gives
\begin{equation}
 E_g(\varepsilon,R):=\int_{T_{\varepsilon,R}}|g|^2\dd V
 =4\pi\int_{\varepsilon<|z-\zeta|<R}
        y^2(|P|^2+|Q|^2)\dd x\dd y.
 \label{eq:energy}
\end{equation}
Put $c=sa+b$ and $\qf_\zeta(a,b)=t^2|a|^2+|c|^2$.
Since $t>0$, $\qf_\zeta$ is positive definite. Equivalently it is
$\|\zeta a+b\|_{\HC}^2$ for the Euclidean complexification norm.

\begin{theorem}[Universal lower bound with optimal leading constant]
\label{thm:lowerenergy}
Let $g$ be any smooth axial Fueter-regular field on the punctured tube;
no growth bound is assumed. Let its periods be $(a,b)$. For
$0<\varepsilon<R<R_0$,
\begin{equation}
 E_g(\varepsilon,R)\geq
 8\int_\varepsilon^R\frac{(t-\rho)^2}{\rho}
 \left(|a|^2+\frac{|sa+b|^2}{t^2+\rho^2}\right)\dd\rho.
 \label{eq:finitelower}
\end{equation}
In particular,
\begin{equation}
 \liminf_{\varepsilon\downarrow0}
 \frac{E_g(\varepsilon,R)}{\log(R/\varepsilon)}
 \geq8\qf_\zeta(a,b).
 \label{eq:lowerlog}
\end{equation}
The constant in \eqref{eq:lowerlog} is sharp for every period pair.
\end{theorem}
\begin{proof}
Parameterize a radius-$\rho$ circle by
$x=s+\rho\cos\varphi$, $y=t+\rho\sin\varphi$.
The definitions of the periods give
\begin{align}
 a&=\frac\rho2\int_0^{2\pi}
       (P\sin\varphi+Q\cos\varphi)\dd\varphi,\\
 c&=\frac\rho2\int_0^{2\pi}
       \big(tP\cos\varphi-(t\sin\varphi+\rho)Q\big)\dd\varphi.
 \label{eq:circleperiods}
\end{align}
For the second equality, add $s\omega_g$ to $\theta_g$ and substitute
$dx=-\rho\sin\varphi\dd\varphi$, $dy=\rho\cos\varphi\dd\varphi$.
The coefficient of $P$ becomes $t\cos\varphi$ and that of $Q$ becomes
$-t\sin\varphi-\rho$.

In the real Hilbert space $L^2([0,2\pi];\R^2)$, the two coefficient
vectors in \eqref{eq:circleperiods} are orthogonal and have squared norms
\[
 \frac{\pi\rho^2}{2},\qquad
 \frac{\pi\rho^2(t^2+\rho^2)}{2},
\]
respectively. Apply Bessel's inequality to each of the four real
quaternionic components and sum. The result is
\begin{equation}
 \int_0^{2\pi}(|P|^2+|Q|^2)\dd\varphi
 \geq\frac{2}{\pi\rho^2}
       \left(|a|^2+\frac{|c|^2}{t^2+\rho^2}\right).
 \label{eq:circlelower}
\end{equation}
In \eqref{eq:energy}, $y^2\geq(t-\rho)^2$ and the planar area element
is $\rho\dd\rho\dd\varphi$. Integration proves
\eqref{eq:finitelower}. The difference between its radial numerator and
$t^2|a|^2+|c|^2$ is $O(\rho)$ as $\rho\downarrow0$, so its contribution
after division by $\rho$ is integrable. This gives \eqref{eq:lowerlog}.
Sharpness follows from the next theorem, applied to $K_\zeta^{a,b}$.
\end{proof}

\begin{theorem}[Exact critical energy and a finite renormalized limit]
\label{thm:exactenergy}
If $g=O(\rho^{-1})$, then for every fixed $0<R<R_0$,
\begin{equation}
 E_g(\varepsilon,R)
 =8\qf_\zeta(a,b)\log(R/\varepsilon)+C_g(R)+o(1)
 \label{eq:exactenergy}
\end{equation}
with a finite real constant $C_g(R)$. Thus subtracting the logarithmic
term gives an actual limit, not merely a bounded remainder.
\end{theorem}
\begin{proof}
The critical normal form reduces the assertion to the local expansion of
$K_\zeta^{a,b}$. Set $d=c/t$. Formula~\eqref{eq:kernelPQ}, with
$x=s+\rho\cos\varphi$ and $y=t+\rho\sin\varphi$, gives uniformly in
$\varphi$
\begin{align}
 P&=\frac{a\sin\varphi+d\cos\varphi}{\pi\rho}
          +O(1+|\log\rho|),\\
 Q&=\frac{a\cos\varphi-d\sin\varphi}{\pi\rho}
          +O(1+|\log\rho|).
 \label{eq:leadingPQ}
\end{align}
The reflected kernel is smooth near $\zeta$, and so is the remainder
in Corollary~\ref{cor:critical}. For the leading coefficient pair,
\[
 |a\sin\varphi+d\cos\varphi|^2
 +|a\cos\varphi-d\sin\varphi|^2=|a|^2+|d|^2.
\]
The cancellation uses the real inner product and does not require $a,d$
to commute. Since $y^2=t^2+O(\rho)$, the leading radial energy density
from \eqref{eq:energy} is
$8(t^2|a|^2+|c|^2)/\rho$.
The difference between the full density and this term is
$O(1+|\log\rho|+\rho\log^2\rho)$, which is integrable at zero.
Its integral therefore has a finite limit. This proves
\eqref{eq:exactenergy}.
\end{proof}

\begin{corollary}[Finite energy annihilates all local periods]
If $E_g(0,R)<\infty$, then $a=b=0$, even without a pointwise growth
hypothesis. More quantitatively, for fixed $R<t$ the lower bound has
the form
\[
 E_g(\varepsilon,R)\geq
 8\qf_\zeta(a,b)\log(R/\varepsilon)-C_{a,b,t,R},
\]
where the constant is independent of $g$ and $\varepsilon$.
\end{corollary}
\begin{proof}
The positivity of $\qf_\zeta$ and
Theorem~\ref{thm:lowerenergy} give the first statement. Integrate the
$O(\rho)$ difference in the proof of that theorem to obtain the second.
\end{proof}

\begin{proposition}[Energy scale of a higher principal part]
\label{prop:poleenergy}
In Theorem~\ref{thm:normalform}, let $k\geq1$ be the largest index with
$C_k\ne0$. Then
\begin{equation}
 E_g(\varepsilon,R)
 =16\pi^2 k\|C_k\|_{\HC}^2\,\varepsilon^{-2k}
       +o(\varepsilon^{-2k}).
 \label{eq:poleenergy}
\end{equation}
Thus the logarithmic defects and rational principal parts have distinct,
completely specified leading $L^2$ scales.
\end{proposition}
\begin{proof}
For a holomorphic stem $F=A+\iota B$ the upper complexified coefficient
pair of its Fueter image is
\[
 P+\iota Q=\frac{2\iota}{y}F'(z)-\frac{2\iota}{y^2}B(z).
\]
The highest pole in $F$ is $(z-\zeta)^{-k}C_k$, so uniformly in the
angle
\[
 |P|^2+|Q|^2
 =\frac{4k^2}{t^2}\|C_k\|^2\rho^{-2k-2}(1+O(\rho)).
\]
Lower poles, the logarithmic defect, and the smooth remainder only give
lower-order terms; the displayed relative estimate may be enlarged by
an $o(1)$ term without affecting the argument. Integrating the leading
term in \eqref{eq:energy} gives
\[
 32\pi^2k^2\|C_k\|^2\int_\varepsilon^R\rho^{-2k-1}\dd\rho
 =16\pi^2k\|C_k\|^2\varepsilon^{-2k}+O(1).
\]
The lower-order contribution is $o(\varepsilon^{-2k})$ by splitting at
a fixed small radius and then letting that radius tend to zero.
\end{proof}

\section{Surface sources and endpoint removability}
\label{sec:source}
The preceding periods also have a distributional interpretation. Let
$\delta_{\Sigma_\zeta}$ denote surface measure as a distribution:
$\langle\delta_{\Sigma_\zeta},\phi\rangle
=\int_{\Sigma_\zeta}\phi\dd S$. For a point
$p=s+tI_p$ on that sphere, $I_p=(p-s)/t$.

\begin{theorem}[The source determined by the two periods]
\label{thm:source}
The locally integrable kernel $K_\zeta^{a,b}$ satisfies
\begin{equation}
 \D K_\zeta^{a,b}
 =2\left(\frac{sa+b}{t}+I_p a\right)\delta_{\Sigma_\zeta}.
 \label{eq:source}
\end{equation}
The same formula holds locally for every critical-growth field with
periods $(a,b)$, after extending it as a locally integrable field across
the sphere. In terms of the normalized Cauchy--Fueter kernel
\[
 E(q)=\frac{\bar q}{2\pi^2|q|^4},\qquad \D E=\delta_0,
\]
one has, for $q\notin\Sigma_\zeta$,
\begin{equation}
 K_\zeta^{a,b}(q)=2\int_{\Sigma_\zeta}
 E(q-p)\left(\frac{sa+b}{t}+I_p a\right)\dd S(p).
 \label{eq:layer}
\end{equation}
\end{theorem}
\begin{proof}
A radius-$\rho$ tube has outward unit normal
$n=\cos\varphi+I\sin\varphi$ and surface element
$y^2\rho\dd\varphi\dd\sigma(I)$. With $d=(sa+b)/t$, the leading term
of \eqref{eq:leadingPQ} obeys
\[
 n(P+IQ)=\frac{d+Ia}{\pi\rho}+O(1+|\log\rho|).
\]
Indeed $(\cos\varphi+I\sin\varphi)(\sin\varphi+I\cos\varphi)=I$
and $(\cos\varphi+I\sin\varphi)(\cos\varphi-I\sin\varphi)=1$.
Integrate by parts on the complement of the tube against a smooth
compactly supported real test function. The inner boundary contributes
with the sign that gives $+nK$ in the distributional derivative. Its
limit is
\[
 \int_{\SSS}2t^2(d+Ia)\phi(s+tI)\dd\sigma(I)
 =2\int_{\Sigma_\zeta}(d+I_pa)\phi(p)\dd S(p).
\]
The error tends to zero because it is multiplied by $\rho$. This proves
\eqref{eq:source}; Corollary~\ref{cor:critical} proves its local version
for a general critical field.

For completeness, the normalization $\D E=\delta_0$ follows by
$\D E=0$ away from zero and the flux integral on a small $3$-sphere:
$nE=1/(2\pi^2r^3)$, whose integral is $1$. Let the right side of
\eqref{eq:layer} be $U$. Convolution of this fundamental solution with
the displayed compact surface source gives $\D U=\D K$ in
distributions. The convolution is locally integrable: in a bounded
four-dimensional region the kernel $|q-p|^{-3}$ is integrable uniformly
for $p$ on a compact surface. Thus $K-U$ is distributionally monogenic,
and applying $\overline\D$ makes each real component harmonic. Weyl's
lemma makes the difference smooth and entire.

Both fields tend to zero at infinity. This is immediate for $U$ from
its integral; for $K$, expand the local stem \eqref{eq:pairedlog} at
infinity. Its Laurent expansion has a constant term and then negative
powers, so its Fueter image is $O(|q|^{-3})$. The parity of the expansion
makes this estimate uniform, including directions near the real axis.
The entire harmonic difference is bounded and hence constant by the
harmonic Liouville theorem; its limit at infinity makes it zero.
\end{proof}

The layer representation makes the relationship with the known
sphere-supported Cauchy kernels in \cite{CPSS14} transparent. It also
shows why two quaternionic parameters occur: the admissible surface
source in this critical axial class has one constant angular mode and
one mode linear in $I_p$.

\begin{theorem}[Sharp $L^2$ removability for a nonreal sphere]
\label{thm:L2}
Let $g$ be Fueter regular on a neighborhood of $\Sigma_\zeta$ with that
sphere removed. If $g$ is locally square integrable across the sphere,
then it extends uniquely to a smooth Fueter-regular function across it.
Axiality is not needed for this removability assertion. If $g$ is axial,
its extension is axial as well. The exponent $2$ cannot be replaced by
any fixed $p$ with $1\leq p<2$: the nonzero fields
$K_\zeta^{a,b}$ are counterexamples.
\end{theorem}
\begin{proof}
Use the distance $\rho$ to the sphere in a small fixed tube. Choose a
cutoff $\chi_\varepsilon$ which is $0$ for
$\rho\leq\varepsilon^2$, $1$ for $\rho\geq\varepsilon$, and affine
in $\log\rho$ between them. It can be smoothed without changing the
bounds
\[
 |\nabla\chi_\varepsilon|
 \leq\frac{C}{\rho\log(1/\varepsilon)},\qquad
 \|\nabla\chi_\varepsilon\|_{L^2}^2
 \leq\frac{C'}{\log(1/\varepsilon)}.
\]
The second estimate uses the codimension-two volume element
$O(\rho\dd\rho\dd\varphi\dd S)$.
Test $\D g=0$ away from the sphere using
$\chi_\varepsilon\phi$, where $\phi$ is smooth and compactly supported
in the tube. The term with $\chi_\varepsilon\nabla\phi$ converges to
the desired distributional derivative. The cutoff error has norm at
most
\[
 C_\phi\|g\|_{L^2}\|\nabla\chi_\varepsilon\|_{L^2}\longrightarrow0.
\]
Thus the locally integrable extension of $g$ satisfies $\D g=0$
distributionally on the full tube. Apply $\overline\D$ and Weyl's lemma
to obtain a smooth harmonic representative, which still satisfies the
first-order equation. Uniqueness follows since two continuous extensions
agree on the dense complement of the sphere. The axial representation
persists by taking limits of the values on each sphere, or by uniqueness
of the smooth extension.

Finally $K_\zeta^{a,b}=O(\rho^{-1})$, so its $p$th power is integrable
for $p<2$ because $\int_0^R\rho^{1-p}\dd\rho<\infty$. It cannot be
removable when $(a,b)\ne(0,0)$, by its periods or by
\eqref{eq:source}. This proves sharpness.
\end{proof}

\begin{remark}[What is and is not a new removability theorem]
The cutoff proof is a standard first-order elliptic capacity argument.
The new emphasis of this paper is not the general mechanism of $L^2$
removability. It is its exact relation, in the axial class, to the two
Fueter-monodromy periods and to the coefficient in
\eqref{eq:exactenergy}. A nonreal two-sphere has codimension two;
these thresholds must not be confused with those for an isolated point
in four dimensions.
\end{remark}

\section{Worked checks and reproducibility}
\label{sec:checks}
The proofs above are analytic. The supplied program
\texttt{checks/verify.py} separately checks the branch-free kernel,
period normalization, the Gram matrix in the energy lower bound, and
the limiting energy coefficient. It does not use symbolic or numerical
success as a substitute for the proofs.

\subsection{Symbolic checks}
SymPy verifies, with symbolic real $x,y,s,t,a,b$, that the expressions
in \eqref{eq:kernelPQ} satisfy both Vekua equations identically. Scalar
$a,b$ suffice here because every displayed identity is real-linear in
each quaternionic component. It also verifies the two diagonal Gram
entries and the zero off-diagonal entry used in
Theorem~\ref{thm:lowerenergy}, and the real-axis limits of
\eqref{eq:explicitcounter}. All seven symbolic checks returned true.

\subsection{Quaternion-valued period checks}
The period test uses
\[
 \zeta=0.3+1.7\iota,\quad
 a=1-2\ii+0.5\jj+3\kk,\quad
 b=-0.7+1.3\ii-2\jj+0.25\kk.
\]
A periodic trapezoidal rule with $4096$ angular points evaluates both
forms on circles about $\zeta$. The maximum absolute error over the
eight real components was
\begin{center}
\begin{tabular}{@{}rr@{}}
\toprule
Circle radius & Maximum absolute period error\\
\midrule
$0.08$ & $4.44\times10^{-15}$\\
$0.31$ & $2.22\times10^{-15}$\\
$0.80$ & $5.55\times10^{-15}$\\
\bottomrule
\end{tabular}
\end{center}
A separate two-hole superposition test, with different quaternionic data
at each hole, had maximum error $1.05\times10^{-14}$. A test circle
surrounding neither hole returned periods below $4.46\times10^{-16}$.
These are floating-point diagnostics, not certified error bounds.

\subsection{Independent source and higher-pole checks}
The additional script \texttt{checks/verify\_sources\_and\_poles.py} evaluates
\eqref{eq:layer} independently by spherical quadrature, implementing the
Hamilton product rather than replacing quaternionic products by scalar
multiplication. For three nonsingular sample points and noncommuting
quaternionic parameters, its maximum component error against the explicit
paired kernel was $2.00\times10^{-15}$. It also checks the leading shell
energy for poles of orders $1$, $2$, and $3$ with a genuinely
complexified quaternionic coefficient. At radius $10^{-4}$ the relative
errors from the predicted coefficients were below $4.33\times10^{-10}$.
The exact input data, quadrature rules, and outputs are in the script and
\texttt{checks/source\_and\_pole\_results.json}.

\subsection{Energy coefficient and renormalization}
For the energy test, take the same $\zeta$ and
\[
 a=1+0.5\ii-0.25\jj+0.75\kk,\qquad
 b=0.2-0.4\ii+0.3\jj-0.1\kk.
\]
The predicted logarithmic coefficient is $8\qf_\zeta(a,b)=46.38$.
The energy per unit logarithmic radial scale is
\[
 S(\rho)=4\pi\rho^2\int_0^{2\pi}
 y^2(|P|^2+|Q|^2)\dd\varphi.
\]
The computation gives
\begin{center}
\begin{tabular}{@{}rrr@{}}
\toprule
$\rho$ & $S(\rho)$ & Relative error from $46.38$\\
\midrule
$10^{-1}$ & $47.21707783$ & $1.80\times10^{-2}$\\
$10^{-2}$ & $46.41552258$ & $7.66\times10^{-4}$\\
$10^{-3}$ & $46.38079719$ & $1.72\times10^{-5}$\\
$10^{-4}$ & $46.38001409$ & $3.04\times10^{-7}$\\
$10^{-5}$ & $46.38000022$ & $4.72\times10^{-9}$\\
\bottomrule
\end{tabular}
\end{center}
With $R=0.45$, the corresponding values of
$E(\varepsilon,R)-46.38\log(R/\varepsilon)$ are $2.33827773$,
$2.35985505$, and $2.36030704$ at
$\varepsilon=10^{-2},10^{-3},10^{-4}$, respectively, consistent with
the proved finite limit.

\subsection{Building and rerunning}
The article is self-contained LaTeX and uses an inline bibliography.
Compile twice with \texttt{pdflatex global\_fueter\_primitives.tex}, or
run \texttt{latexmk -pdf global\_fueter\_primitives.tex}. To reproduce
the diagnostic output, install NumPy, SymPy and SciPy and run
\begin{quote}\ttfamily\small
python checks/verify.py --output checks/verification\_results.json
\end{quote}
The archive includes the exact JSON output, a requirements file, the
literature-search log, and the supplied manuscript unchanged in a
separate source-context directory. No theorem here has been checked by
a proof assistant.

\section{Consequences and precise limits of the results}
\label{sec:conclusion}
The local inverse problem has a particularly simple global obstruction
because its kernel is affine. Passing to $C=B/y$ and $V=A-xB/y$
turns the two unknown affine coefficients into two ordinary closed
one-forms. Reflection then removes the apparent obstruction from holes
that meet the real axis and doubles neither the number of independent
parameters nor the required correction family.

Three conclusions are worth keeping distinct. The global theorem
classifies the \emph{range} of the Fueter map on the same domain; it is
not a statement that arbitrary nonaxial monogenic functions are Fueter
images. The approximation theorem classifies which fixed defects must
be added to rational Fueter images; it does not change the meaning of
slice regularity or allow multivalued rational approximants. The
singularity theorem classifies finite power growth near a nonreal
sphere; it does not assume that all isolated spherical singularities
have finite order.

Within those limits, the results answer the opening questions
completely. The obstruction is explicitly computable, every allowed
obstruction occurs, its quaternionic dimension is known, its normal
form is unique, and its critical energy cost is sharp. The local
existence and the elementary kernels remain part of the established
Fueter theory; the reflection-sensitive range statement and the
quantitative classification are the research claims to compare with
further literature.

\appendix
\section{An implementation-oriented reconstruction procedure}
Given smooth axial data $P,Q$ on a finite-type parameter domain, first
check the two Vekua residuals. Form the coefficient pairs of
$\omega_g$ and $\theta_g$ directly; no primitive is needed at this
stage. Evaluate their periods only on loops about one upper hole in
each conjugate pair. Periods about fixed holes vanish analytically;
numerical nonzero values there detect data or quadrature error.

Let the computed exact-data period vector be $(a_j,b_j)$. Subtract
$\sum_jK_{\zeta_j}^{a_j,b_j}$, integrate the two resulting forms along
paths from a basepoint, and set $f(q)=qC+V$. In an exact computation,
path independence follows from Theorem~\ref{thm:range}. In a numerical
implementation, differences between integrals over distinct paths
provide a consistency diagnostic. Inequality~\eqref{eq:inversebound}
controls amplification by path length, while
\eqref{eq:nondensity} prevents an algorithm from silently treating
nonzero monodromy as ordinary approximation error.

Near a sphere, avoid direct subtraction of nearly equal terms in the
paired formula when $y$ is close to zero. The local intrinsic stem
\eqref{eq:pairedlog}, its real-centered expansion, or the Cauchy-layer
formula is a more stable representation there. Near the singular
sphere, use the explicit leading part \eqref{eq:leadingPQ} plus the
remaining logarithmic terms. These are numerical representation choices,
not additional analytic assumptions.

\begin{thebibliography}{10}
\addcontentsline{toc}{section}{References}

\bibitem{Attachment}
\emph{Classical Complex Theorems over the Quaternions: Slice-regular and
Cauchy--Fueter analogues, with proofs}.
Supplied expository manuscript, dated 20 September 2026,
file \texttt{quaternionic\_analysis(1).tex}, especially Section~13,
``Fueter's mapping theorem and a constructive local inverse.''
Included unchanged in the source-context directory of this archive.

\bibitem{GP11}
R.~Ghiloni and A.~Perotti,
\emph{Slice regular functions on real alternative algebras},
Advances in Mathematics \textbf{226} (2011), 1662--1691.
\href{https://arxiv.org/abs/1008.4318}{arXiv:1008.4318}.

\bibitem{CSS11}
F.~Colombo, I.~Sabadini, and F.~Sommen,
\emph{The inverse Fueter mapping theorem},
Communications on Pure and Applied Analysis \textbf{10} (2011), 1165--1181.
\href{https://doi.org/10.3934/cpaa.2011.10.1165}{doi:10.3934/cpaa.2011.10.1165}.

\bibitem{CSS13}
F.~Colombo, I.~Sabadini, and F.~Sommen,
\emph{The inverse Fueter mapping theorem in integral form using spherical
monogenics}, Israel Journal of Mathematics \textbf{194} (2013), 485--505.
\href{https://doi.org/10.1007/s11856-012-0090-4}{doi:10.1007/s11856-012-0090-4}.

\bibitem{CPSS14}
F.~Colombo, D.~Pe\~na Pe\~na, I.~Sabadini, and F.~Sommen,
\emph{A new integral formula for the inverse Fueter mapping theorem},
Journal of Mathematical Analysis and Applications \textbf{417} (2014),
112--122.
\href{https://doi.org/10.1016/j.jmaa.2014.03.016}{doi:10.1016/j.jmaa.2014.03.016};
\href{https://arxiv.org/abs/1302.0685}{arXiv:1302.0685}.

\bibitem{Perotti24}
A.~Perotti,
\emph{A local Cauchy integral formula for slice-regular functions},
Computational Methods and Function Theory \textbf{24} (2024), 185--203;
online publication 30 May 2023.
\href{https://doi.org/10.1007/s40315-023-00485-5}{doi:10.1007/s40315-023-00485-5}.

\bibitem{DMP26}
A.~De Martino and S.~Pinton,
\emph{A variation of the inverse Fueter theorem and the generalized
polyanalytic Cauchy--Kovalevskaya extension of order 2}, preprint,
7 August 2026.
\href{https://arxiv.org/abs/2608.07711}{arXiv:2608.07711}.

\end{thebibliography}
\end{document}
